% INSERTION: r1.pdf, p. 11, immediately after the paragraph below the
% two decay-constant sum rules ending "presented in Appendix C."
% Place before "Substituting the residues obtained from the two-point sum rules..."
% Aliev:2025mixing is the key in the local bibliography. Check the Prism key.
% This summarizes the earlier angle method; it does not report an input update.

For completeness, we briefly summarize the determination of the mixing angle
in Ref.~\cite{Aliev:2025mixing}. The two basis currents have the same mass
dimension, with the tensor current normalized by $1/(m_Q+m_s)$.\footnote{This
factor was inadvertently omitted from the printed definition
of the tensor current in Ref.~\cite{Aliev:2025mixing}, but was included in
the calculations underlying that work.}
In the two-state description adopted in that analysis, the angle is determined
by imposing a vanishing off-diagonal spin-one correlator in the rotated basis.
With the current convention used here, the rotation is
\begin{equation}
 R(\theta)=
 \begin{pmatrix}
  \sin\theta & \cos\theta\\
  \cos\theta & -\sin\theta
 \end{pmatrix}.
 \label{eq:r1-angle-rotation}
\end{equation}
Writing the Borel-transformed and continuum-subtracted correlator matrix as
$\widehat{\boldsymbol\Pi}$, the condition is
\begin{equation}
 R(\theta)\widehat{\boldsymbol\Pi}R^T(\theta)
 =\begin{pmatrix}
   \widehat\Pi_{11}&0\\0&\widehat\Pi_{22}
  \end{pmatrix},
 \qquad
 \widehat{\boldsymbol\Pi}=
 \begin{pmatrix}
  \widehat\Pi_{AA}&\widehat\Pi_{AB}\\
  \widehat\Pi_{AB}&\widehat\Pi_{BB}
 \end{pmatrix},
 \label{eq:r1-angle-diagonalization}
\end{equation}
where the matrix entries in the angle sum rule are evaluated with a common
Borel parameter and continuum-subtraction prescription. The vanishing
off-diagonal entry gives
\begin{equation}
 \frac{1}{2}\sin(2\theta)
 \bigl(\widehat\Pi_{AA}-\widehat\Pi_{BB}\bigr)
 +\cos(2\theta)\widehat\Pi_{AB}=0.
 \label{eq:r1-angle-offdiagonal}
\end{equation}
This condition does not require the two physical diagonal correlators to be
equal. Indeed, defining
$\Delta=\widehat\Pi_{11}-\widehat\Pi_{22}$, the inverse rotation yields
\begin{align}
 \widehat\Pi_{AA}-\widehat\Pi_{BB}&=-\cos(2\theta)\Delta,
 \nonumber\\
 2\widehat\Pi_{AB}&=\sin(2\theta)\Delta.
 \label{eq:r1-angle-inverse}
\end{align}
For nondegenerate correlators, their common factor $\Delta$ cancels in
the ratio, giving
\begin{equation}
 \tan(2\theta)=
 -\frac{2\widehat\Pi_{AB}(M^2,s_0)}
 {\widehat\Pi_{AA}(M^2,s_0)-\widehat\Pi_{BB}(M^2,s_0)}.
 \label{eq:r1-angle-sumrule}
\end{equation}
Thus, different pole masses and residues are compatible with this relation.
The mixing angles obtained in Ref.~\cite{Aliev:2025mixing} are used as external
inputs in the present analysis and are quoted in the numerical section below.
